\section{Introducción: por qué emprender en robótica «no es romántico»}

Esta sección fija primero la manera de leer el episodio: no conviene tomarlo como una noticia de chismes sobre «quién dejó qué empresa», ni solo como publicidad corporativa de Xinghaitu (星海图). Una lectura más valiosa es verlo como el testimonio oral de cómo un emprendedor traslada su formación en conducción autónoma a la inteligencia corporizada.

En apariencia, esta entrevista trata de la trayectoria personal de Gao Jiyang (高继扬), fundador de Xinghaitu; en realidad es una clase sobre la industria acerca de «cómo crece una empresa de IA del mundo físico». Al inicio, el conductor plantea una duda: ¿por qué en el sector chino de la inteligencia corporizada casi no aparecen personas con el marcado romanticismo tecnológico que se ve en los emprendimientos de modelos grandes? La respuesta de Gao Jiyang atraviesa toda la conversación: la cadena de la robótica es larguísima, los ciclos son larguísimos y el equipo tiene, por naturaleza, que meter la cabeza en la tierra.

Aquí «en la tierra» no es una expresión negativa, sino las restricciones reales del mundo físico: el equipo completo, la cadena de suministro, la recolección de datos, las instalaciones del cliente, la latencia en el dispositivo, el mantenimiento, los canales de distribución, el financiamiento y el crecimiento de la organización. Los modelos grandes o las aplicaciones de software pueden arrancar a partir del modelo y de la difusión, pero a una empresa de robótica le resulta muy difícil esquivar el hardware y el valor para el cliente. Por eso, estas notas organizan el episodio como un «campo de entrenamiento industrial de la IA del mundo físico»: Waymo le enseñó ingeniería de sistemas, Momenta le enseñó la entrega en producción en masa, y Xinghaitu recombina el equipo completo, los datos, los modelos y el ciclo cerrado con el cliente.

\begin{importantbox}{Tesis central del episodio}
La inteligencia corporizada no es «una empresa de algoritmos con una carcasa de robot», sino una cadena larga que va del equipo completo, la cadena de suministro, el sistema de datos, la AI infra y los modelos algorítmicos hasta los canales de distribución y el valor para el cliente. A corto plazo los algoritmos pueden difundirse con rapidez; a largo plazo, la barrera competitiva tiene más probabilidades de surgir del ciclo cerrado con el mundo real.
\end{importantbox}

\begin{knowledgebox}{Nota sobre la estrategia visual}
Este video es una toma fija de entrevista, sin slides, pizarrón ni demostraciones de producto. Conforme al estándar para notas de pódcast, el cuerpo del texto no repite fotogramas de las personas; la portada sirve para identificar la fuente, y en el cuerpo se usan diagramas conceptuales, como el modelo de negocio de la conducción autónoma, el volante de datos y el sistema dual del cerebro robótico, para explicar el contenido de la entrevista.
\end{knowledgebox}

\subsection{Resumen de la sección}

Este episodio no es una simple historia de fundador, sino un caso de cómo la experiencia en conducción autónoma se traslada a la inteligencia corporizada. Más adelante se abordarán, en orden, la metodología de Gao Jiyang, la comparación entre Waymo y Momenta, las decisiones de Xinghaitu sobre el equipo completo y la cadena de suministro, la Data Recipe, el cerebro robótico VLM/VLA y las disyuntivas organizacionales detrás de «la salida de Xu Huazhe (许华哲)».
